\section{Knowledge organization and lifecycle management}
\label{sec:organization}

Acquisition decides what enters the system; organization decides what the
system is. In this platform organization is not a folder convention that a
human maintains by hand, nor an emergent property of embedding similarity.
It is a governed process with three invariants: content decides where a
document belongs, the disk, the tree registry and the library metadata must
agree about every document, and every operation that removes or relocates
knowledge is preceded by a survey and followed by a verification. This
chapter describes how those invariants are implemented, which of them are
enforced by the interface rather than by convention, and where the current
implementation still allows the three stores to diverge.

\subsection{Content-routed placement}
\label{sec:org-placement}

Placement reads the parsed document, not its filename or its folder of
origin. For each arriving document the pipeline selects the library whose
stated scope covers the content, or proposes a new library when no existing
scope does; it then writes the document, derives a description and a tag
set, and runs the same final check used by acquisition before the document
counts as placed. Two rules keep the catalog usable as a retrieval
instrument. First, an existing library is preferred over a new one: the
catalog stays shallow, and the number of libraries an agent must inspect
before searching stays small. Second, a document is placed once; re-routing
is an explicit lifecycle operation, not a side effect of a later ingest.

The library tree supports one level of content grouping in which a parent
library contains child libraries, and the metadata of a child is addressable
through the parent. The implementation carries a caveat that we report
rather than hide: the vector segments of child documents are stored in the
parent library's collection, so a search that targets a child library
identifier returns no candidates. The shipped operational guidance is
therefore to search the parent collection and to use the child libraries for
navigation and scoping rather than as search endpoints. This is a limitation
of the current storage mapping, not of the routing rule, and it is the kind
of coupling between an organizational decision and a retrieval path that a
system paper has to expose.

\subsection{Catalog and metadata registries}
\label{sec:org-registries}

Two registries describe the same collection, and Table~\ref{tab:datamodel}
lists their fields: the global tree (\texttt{.tree-fs.yml}) is the
authoritative structure, and each library carries a
\texttt{.knowledge-base.yml} that indexes its documents. Documents themselves
are Markdown files with an accompanying image directory.

Agents read this catalog through a progressive-disclosure interface: the
library listing returns names, descriptions and document counts; the
document listing returns metadata without bodies; only an explicit read
returns content, and a read is attributed to a specific document. This design
choice is what makes the catalog-first retrieval lane viable: a caller can
narrow a search to a handful of candidate documents without ever loading a
body, which is why the catalog lane completes in 5--6 turns and 54--91~s
(Table~\ref{tab:lanes}).

Correctness of the catalog is defined as three-source consistency: for every
document, the file on disk, its entry in the tree registry, and its record in
the library metadata must agree on existence, path, size and index state. The
platform verifies this agreement rather than assuming it, and the audit is
productive: in the acceptance campaign a renamed document kept its old vector
segments, so the recorded chunk count read 18 where 9 chunks existed, and the
document matched its own former path at 0.9738 similarity, which is high
enough to be reported as a near-duplicate pair. Only an explicit force
re-index repaired it (18~$\rightarrow$~9). The design lesson is that metadata
consistency and index consistency are two different obligations: the first
can be satisfied by a single write operation, the second cannot, because the
vector store is updated by a separate call.

\subsection{Tag taxonomy and governance}
\label{sec:org-tags}

Tags are the coarse, human-legible axis of the collection. A tag registry
holds the vocabulary; each document's tags are normalized before they are
written, filtered against a blacklist of structural strings (section
headings, generic words such as ``introduction''), and re-checked against the
document body so that a tag must actually occur in the content it labels.
Retrieval by tag is a first-class operation, which makes tag quality a
retrieval concern rather than a cosmetic one.

Governance is a maintenance loop rather than a schema. The registry is
curated, not centralized: no taxonomy is imposed on authors, and the
vocabulary grows with the corpus. The audit therefore treats the registry as
a health surface: it reports tags with zero referencing documents, which
accumulate from test documents, abandoned naming experiments and chapter
headings, and it removes them only when instructed. In the audited campaign
165 documents passed the tag gate with no failures, and the live registry
held 316 tags. We do not claim a controlled vocabulary; we claim that the
vocabulary is inspectable and that orphaned entries are detectable.

\subsection{Duplicate detection and collection hygiene}
\label{sec:org-dedup}

Acquisition refuses exact duplicates by content hash and flags near
duplicates by vector similarity; at the collection level the question is
different, because the same content can legitimately arrive through two
ingest paths or into two libraries. The hygiene procedure is
survey-then-act: enumerate candidate duplicate groups at a similarity
threshold (0.90, with a manual fallback channel at 0.85), inspect each group,
and delete nothing automatically.

The audited run illustrates why the survey step matters more than the
threshold. On the 322-document corpus the detector reported 35 candidate
groups with similarities between 0.90 and 0.975; every one of them turned out
to be a pair of adjacent parts of the same long paper rather than a true
duplicate, and no deletion was performed. A policy that deleted on detection
would have destroyed 35 legitimate document parts. The same run confirmed
zero true duplicates, which is the expected outcome when acquisition gates
run before storage. Cross-library duplication is detected but not merged
automatically: merging two libraries is a lifecycle decision with provenance
consequences, so it stays with the operator.

\subsection{Library lifecycle operations}
\label{sec:org-lifecycle}

Table~\ref{tab:org-lifecycle} lists the operations the platform exposes, the
interface that performs them, and the consistency obligation each one
carries. All of them are available to agents through the same tool surface
that the user interface uses, which is deliberate: an agent that may read a
collection should be able to maintain it under the same rules a human
operator follows.

\begin{table}[t]
\centering
\small
\caption{Library and document lifecycle operations, their consistency
obligation, and the state observed in the acceptance campaign. ``Audited''
means the three stores were compared after the operation.}
\label{tab:org-lifecycle}
\begin{tabular}{@{}p{0.22\textwidth}p{0.24\textwidth}p{0.20\textwidth}p{0.24\textwidth}@{}}
\toprule
Operation & Interface & Obligation & Observed \\
\midrule
Create library & library create & tree + metadata + directory & audited \\
Rename or move document & document move & path updated in all three stores & audited; stale segments until force re-index \\
Update description or tags & document update & metadata rewritten, body re-checked & audited (165/165) \\
Update content & document update & new hash, re-index required & re-index not implicit \\
Delete document & document delete & removal from all three stores & audited (1~$\rightarrow$~0) \\
Batch delete & batch delete & per-document verification & 1/1 in campaign \\
Move library & library move & subtree and metadata move together & audited (0~$\rightarrow$~1) \\
Merge libraries & library merge & provenance preserved & demonstrated, not controlled \\
Split library & library split & no part orphaned & analysed, deliberately not executed \\
Rebuild index & re-index (task id) & segment count matches content & 18~$\rightarrow$~9 after force \\
\bottomrule
\end{tabular}
\end{table}

Two properties of this table carry the argument of the chapter. The first is
that re-indexing is not implicit. Creating, moving or rewriting a document
does not update the vector store; an explicit index operation does, and the
audit treats a missing or stale index as an incomplete operation rather than
as a warning. This is the single rule we would carry into any successor
system: an index that is updated by side effect is an index that silently
disagrees with the truth. Keeping an approximate index coherent with a
mutating document store is itself an active research problem
\citep{xu-spfresh-2023}; the contribution here is not a new maintenance
algorithm but a contract under which disagreement becomes detectable. The
second is the honest boundary of the evidence.
Move, rename, update, delete and re-index were exercised with real tool calls
and audited afterwards; merge and split are implemented and were demonstrated
on live libraries, but they have no controlled evaluation, and the
sub-library hierarchy was analysed and deliberately not executed pending an
operator decision. No concurrent organizational operation was tested. We
prefer to state this than to imply that a lifecycle operation is verified
because it has an interface.

\subsection{The document graph}
\label{sec:org-graph}

The platform maintains a property graph over the collection with document,
library, tag, folder and entity nodes and membership, tagging, relatedness,
identity and containment edges. Edges are derived from three signals:
co-membership of tags, vector similarity above a threshold, and the verdicts
produced by content verification during retrieval. Adding one document to a
library added four edges in the audited run (4{,}675~$\rightarrow$~4{,}679),
one per assigned tag; a full build over the 100-paper corpus produced 658
nodes covering 333 documents and 316 tags, and the 165-document intermediate
state held 335 nodes and 1{,}671 edges.

The graph has a deliberately restricted role. It is a navigation instrument:
it expands a candidate set, surfaces documents that share context across
libraries, and identifies central documents in a collection. It is never
evidence. Every candidate a graph traversal proposes still has to pass the
same content-verification gate as a candidate proposed by a dense retriever,
and the graph's own contribution to retrieval is measured separately: as a
standalone retriever it reached Hit@5 of 0.5625 on the industrial corpus and
0.2 on the second corpus, well below the lexical and dense channels
(Table~\ref{tab:retrieval}). We report this because it is the honest status
of graph-based organization in this system: useful for discovery, weak as a
ranker.

Two implementation gaps belong here as well. Graph construction is not yet
guaranteed by the ingestion path: a document ingested through one agent
harness received no graph node, so a library reported two of three
documents, which is a silent incompleteness of the kind the three-source
audit is designed to catch elsewhere. And batch graph builds have been
observed to leave a task record in the running state after the work
completed, so task status cannot be used as completion evidence; the graph
itself must be queried.

\subsection{Operational knowledge as a first-class collection}
\label{sec:org-experience}

Process documents describe what the organization intends; operational
lessons describe what actually happened. The platform stores the second kind
as a first-class collection with the same addressing, indexing and retrieval
interface as ordinary documents, plus two differences that reflect the
epistemic status of a lesson.

First, an experience carries a credibility tier, and the tier controls
whether it may be shown (Table~\ref{tab:exp-tiers}). Knowledge that is a
claim about practice is held to a stricter bar than a document: the
experience channel's similarity threshold is 0.55 against the document
channel's 0.35, and an unvetted or disputed entry is capped even when it
scores well. Second, entries are drafted before they are served. Extraction
produces candidates from ingested content or from recurring question-answer
pairs; automated distillation runs as a separate agent process on a schedule,
guarded by a circuit breaker (three consecutive failures suspend it for 24
hours) and a heuristic fallback whose output is marked low-confidence.
Nothing reaches the served collection without an explicit approval step,
which is the same fail-closed posture the retrieval gate applies to
citations.

\begin{table}[t]
\centering
\small
\caption{Experience credibility tiers and modifiers. Tiers combine retrieval
score, content-verification score, ratings and review count; modifiers cap
the tier rather than promote it.}
\label{tab:exp-tiers}
\begin{tabular}{@{}llp{0.46\textwidth}@{}}
\toprule
Tier & Condition & Effect \\
\midrule
P0 strong & sim $\ge 0.65$ $\wedge$ content $\ge 6$ $\wedge$ rating $\ge 4$ $\wedge$ reviews $\ge 1$ & may be presented first \\
P1 reference & sim $\ge 0.45$ $\wedge$ content $\ge 4$ & presented with credibility noted \\
P2 weak & sim $\ge 0.35$ $\wedge$ content $\ge 3$ & suppressed unless expanded \\
discard & content verification fails or sim $< 0.35$ & never presented \\
\addlinespace
disputed & reviews $\ge 3$ $\wedge$ rating $< 2.0$ & capped at P2 \\
unvetted & 0 reviews $\wedge$ 0 applications & capped at P1 \\
auto-extracted & not human-reviewed & capped at P1 until applied twice and rated $\ge 4$ \\
\bottomrule
\end{tabular}
\end{table}

The evidence for this subsystem is deliberately presented as case evidence.
One extraction campaign produced thirteen candidate drafts; one end-to-end
scenario sedimented a single experience, which a later global experience
search returned first with a local similarity of 0.711; and the approval path
was never exercised, because the draft pool was empty when it was tested. A
single scenario and thirteen candidates cannot support an effectiveness
claim, and we do not make one. What the campaign does establish is that
operational knowledge travels through the same pipeline as documents ---
addressed, indexed, gated and auditable --- rather than living in a wiki or a
prompt.

\subsection{Integrity verification as routine}
\label{sec:org-integrity}

Every property claimed in this chapter is checked by a verification pass that
runs against the live collection: three-source consistency per document,
document integrity, parse quality, index coverage with optional repair, tag
health, experience health, and a scorecard that summarizes the findings. The
pass is read-only by default; repairs require an explicit instruction, and
the report is produced before any repair. This is a deliberate ordering: an
audit that repairs as it goes cannot report what it found.

The pass is not ceremonial. It is how the two index-integrity incidents of
this deployment were found --- a silently lost collection and a stale index
--- and how the reorganization campaign was certified after the fact
(322/322 metadata checks, 165/165 description checks, 33/33 sampled
consistency checks, 15/15 retrieval regression probes with no degradation).
In the same campaign the audit also surfaced the failures reported above: the
rename-time chunk duplication, the missing graph node, and a task-status
record that outlived its task. A collection whose health is a command is a
collection whose health can be argued about with evidence.

\subsection{What organization buys, and what it does not}
\label{sec:org-summary}

Organization buys candidate quality and auditability. Because the catalog is
curated and consistent, an agent can select a library and a document set
without loading bodies, which is the mechanism behind the catalog-first lane
and behind cross-library discovery in an environment with nineteen libraries.
Because every placement, tag, description and deletion left a record, the
collection can be re-verified after the fact; the numbers in this chapter are
not survey estimates but audit outputs.

Three things organization does not do, and the paper does not claim. It does
not repair a bad parse: a document whose parse lost its tables cannot be
recovered by placing it well. It does not prevent divergence by
construction: the write path can leave the vector index behind, and only an
audit and an explicit re-index restore agreement. And it does not scale
without limits: the largest collection we report is 322 documents in five
libraries, the merge and split operations are unevaluated, and no
organizational operation was tested under concurrency. These are engineering
boundaries, and the next chapter describes the retrieval contract that is
built on top of them.
